\subsection{The lower bound by finite weighted flow}\label{sec:rate-lower}

We prove, with constants depending only on the fixed Bell functional,
\begin{equation}\label{eq:lb}
 S-S_d\ge c_* e^{-C_*d}\qquad(d\ge1).
\end{equation}
The proof uses neither a dimension-preserving conversion to a chain nor
identification of vectors in different joint spectral cells. There can be
quadratically many such cells. We count only the distinct local spectral
labels, and control all nonnegative cell weights simultaneously.

\subsubsection*{Critical data and the original-dimension interface}

Write
\[
 \delta(x,u)=xu+(x-u)/2-1,\quad b(t)=\tfrac12\sqrt{1-t^2},
 \quad w=b^2.
\]
The critical-storage construction of Section~\ref{sec:setup} gives a
concave $3/2$-Lipschitz function $g$ on $[-1,1]$, positive in the open
interval, such that
\begin{equation}\label{eq:wf-feasibility}
 \frac{w(x)}{g(x)}+g(u)\le S-\delta(x,u),\qquad
 g(t)g(-t)\ge w(t).
\end{equation}
Initially the quotient is taken at interior points. To make the inherited
input precise, for $q>S$ take the infimum $g_q$ of terminal Schur pivots
over finite Jacobi histories. The bound $J\preceq SI$ gives pivots at least
$q-S$; extending a history gives the first inequality at level $q$.
Each terminal-pivot function has slope $1/2-x$ in the terminal label, so
their infimum is concave and $3/2$-Lipschitz. Monotonicity in $q$ and the
common modulus give a uniform limit $g$ as $q\downarrow S$.
Feasibility with a fixed target label bounds $w(t)/g_q(t)$ uniformly for
each interior $t$, proving $g(t)>0$ there. Concatenating a history ending
at $t$ with the reflected reverse of a history ending at $-t$ gives two
central Schur pivots coupled by $b(t)$. The resulting positive
$2\times2$ matrix gives the second inequality in
\eqref{eq:wf-feasibility} after taking the two infima and the limit.
This is the critical-storage input, not a consequence of the old
finite-dimensional rate argument.
For a related treatment of canonical storage and positive welds for
reversible Bell functionals, see Mghirbi \cite{MghirbiDynamics}.
The argument below supplies the unrestricted lower bound directly;
no such bound is imported from that work.

Here is an elementary endpoint check using only $S<1/3$.
One-edge histories give $g(1),g(-1)\le S$. If $g(1)=0$, Lipschitz
continuity and reflection gluing imply, for $-1<t<1$,
\[
 g(t)\le\tfrac32(1-t),\qquad
 g(-t)\ge\frac{1-t^2}{4g(t)}\ge\frac{1+t}{6}.
\]
Letting $t\uparrow1$ contradicts $g(-1)\le S<1/3$. Reflecting proves
$g(-1)>0$ also. Thus $m:=\min_{[-1,1]}g>0$, and all quotients below
extend continuously to the endpoints.

Set
\[
 p=w/g,\quad \alpha(t)=b(t)\sqrt{g(-t)/g(t)},\quad
 \beta(t)=\alpha(-t),\quad
 \phi(x,u)=S-\delta(x,u)-\alpha(x)-\beta(u).
\]
Feasibility and its reflected version, followed by Cauchy--Schwarz, give
$\phi\ge0$. In detail, the two vectors
$(\sqrt{p(x)},\sqrt{g(u)})$ and
$(\sqrt{g(-x)},\sqrt{p(-u)})$ have inner product
$\alpha(x)+\beta(u)$ and squared norms at most $S-\delta(x,u)$.
At $x=\pm1$ they cannot be proportional: the first coordinate of the
first vector vanishes, whereas both $g(u)$ and $g(-x)$ are positive.
Consequently $\phi>0$ on those two boundary edges; reflection treats the
other two. Compactness supplies a symmetric interval
$K=[-a,a]\subset(-1,1)$ containing every zero strictly inside $K^2$,
and a constant $a_0>0$ with
\begin{equation}\label{eq:wf-collar}
 \inf_{(x,u)\notin K^2}\phi(x,u)\ge a_0.
\end{equation}
For example, minimize on the closed boundary collar
$\max(|x|,|u|)\ge a$. No minimum on a nonclosed set is needed.

First reduce arbitrary binary effects and mixed states without dimension
loss. At fixed state the Bell value is affine in each effect; successively
maximize each coordinate at a projection, then choose a top eigenvector
of the resulting Bell operator. The value cannot decrease and neither
local dimension changes. It therefore suffices to prove the lower bound
for projective pure-state strategies. For such measurements on the
original tensor product, put
\[
 X=A_1+A_2-I,\quad Y=A_2-A_1,\quad
 U=B_1+B_2-I,\quad V=B_2-B_1,
\]
and $W_A=Y(B_3-I/2)$, $W_B=(A_3-I/2)V$. Direct multiplication, using
idempotence of the measurement projections, gives
\[
 W_A^*=W_A,\quad W_AX=-XW_A,\quad W_A^2=b(X)^2,
\]
and the analogous identities for $W_B,U$. In particular $W_A$ is not
assumed to commute with $U$. The operators
$Q_A=\alpha(X)-W_A$ and $T_A=\alpha(-X)+W_A$ are self-adjoint and
satisfy $Q_AT_A=0$: use $W_Af(X)=f(-X)W_A$ and
$\alpha(t)\alpha(-t)=b(t)^2$. They therefore commute. Their sum is
nonnegative, so their joint spectral values show that both factors are
nonnegative. Likewise $Q_B=\beta(U)-W_B\ge0$. Expanding gives
\begin{equation}\label{eq:wf-weld}
 SI-\B=\phi(X,U)+Q_A+Q_B.
\end{equation}
For a unit vector of deficit $\eps=S-\langle\B\rangle$ and its joint
spectral probability measure $\mu$, positivity and bounded norms imply
\begin{equation}\label{eq:wf-residual}
 \int\phi\,d\mu\le\eps,\qquad
 \|Q_A\psi\|+\|Q_B\psi\|\le C\sqrt\eps,\qquad
 \mu((K^2)^c)\le a_0^{-1}\eps.
\end{equation}
Indeed $\|Q_A\psi\|^2\le\|Q_A\|\langle Q_A\rangle$, with a fixed,
dimension-independent norm bound. This derives the approximate response
estimate from the positive weld, not from an exact-equality statement.

If the original local dimensions are at most $d$, each of $X,U$ has at most $d$ distinct
eigenvalues, regardless of spectral multiplicities.

\subsubsection*{A quantitative monotone predictor}

Put $C_0(x)=S+1-x/2-p(x)$ and let $H$ be its greatest convex minorant
on the full interval $[-1,1]$. There is a fixed $M\ge1$ such that
$h=H'$ is nondecreasing and $M$-Lipschitz in the interior, with a
continuous constant extension past either endpoint.
For completeness, for smooth positive concave approximants to $g$,
\[
 p''=\frac{w''}{g}-\frac{2w'g'}{g^2}
       +\frac{2w(g')^2}{g^3}-\frac{wg''}{g^2}
       \ge-\frac1m-\frac{4L}{m^2},
\]
where $L$ bounds the slopes and the approximants are at least $m/2$.
Concave affine extensions followed by mollification and uniform limits
show that $C_0$ is semiconcave. One may take
$M=1+1/m+4L/m^2$. On every open gap where $H<C_0$, maximality makes
$H$ affine. At an interior contact,
\[
 C'_{0,-}\le H'_-\le H'_+\le C'_{0,+},
\]
whereas semiconcavity reverses the outside inequality. All four values
coincide. Adding the semiconcavity inequalities at contact points $x<y$
gives $0\le H'(y)-H'(x)\le M(y-x)$. The affine gaps extend this estimate
to arbitrary points, establishing the assertion.

The affine minorant at $\min_z(C_0(z)-uz)$ lies below $H$, so
\[
 g(u)-u/2\le\min_z(C_0(z)-uz)=\min_z(H(z)-uz).
\]
Write $D=S-\delta(x,u)$, $L_1=D-p(x)-g(u)$ and
$L_2=D-g(-x)-p(-u)$. Both are nonnegative, and
\[
 \alpha(x)+\beta(u)\le\sqrt{(D-L_1)(D-L_2)}
 \le D-(L_1+L_2)/2.
\]
Consequently
\begin{equation}\label{eq:wf-envelope-gap}
 H(x)-ux-\min_z(H(z)-uz)\le L_1\le2\phi(x,u).
\end{equation}
Choose $\rho>0$ with every $x\in K$ at least $\rho$ from the interval
endpoints. For $F(z)=H(z)-uz$, the admissible step
$z=x-\rho\operatorname{sign}F'(x)$ shows that a gap smaller than
$M\rho^2/2$ implies $|F'(x)|<M\rho$. The full step $x-F'(x)/M$ is
then admissible; the smoothness inequality gives a decrease at least
$|F'(x)|^2/(2M)$. Thus, for every $s$ below a fixed sufficiently small
positive threshold,
\begin{equation}\label{eq:wf-predictor}
 x,u\in K,\ \phi(x,u)\le s
 \quad\Longrightarrow\quad |u-h(x)|\le2\sqrt{Ms}.
\end{equation}

\subsubsection*{Scalar balance without identification of joint cells}

Let $p_X,p_U$ be the original spectral marginals, and put
$r(t)=\sqrt{g(-t)/g(t)}$, $R(t)=r(t)^2$. On $K$, $R$ is Lipschitz,
with $0<R_{\min}\le R\le R_{\max}<\infty$.
The operator $b(X)^{-1}W_A$ is an isometry on the $X$-corridor and
reflects its spectral projections. Projecting the response residual in
\eqref{eq:wf-residual}, dividing by $b(t)\ge b_0>0$, and using the
reverse triangle inequality yields
\[
 \sum_{t\in K}
 \left(\sqrt{p_X(-t)}-r(t)\sqrt{p_X(t)}\right)^2\le C\eps.
\]
Missing spectral atoms have weight zero. Cauchy--Schwarz converts this
to the total $\ell^1$ bounds
\[
 \|\mathcal Fp_X-Rp_X\|_{1,K}\le C_1\sqrt\eps,\qquad
 \|\mathcal Fp_U-R^{-1}p_U\|_{1,K}\le C_2\sqrt\eps,
\]
where $\mathcal F$ reflects a measure. These inequalities use genuine
projected response vectors but not their localization in joint cells.

Average measures only:
$\bar\mu=(\mu+J_*\mu)/2$, $J(x,u)=(-u,-x)$. Since
$\phi\circ J=\phi$, the defect bound is unchanged. Its marginals
$p=(p_X+\mathcal Fp_U)/2$, $q=(p_U+\mathcal Fp_X)/2$ satisfy
\[
 q-Rp=\tfrac12(\mathcal Fp_X-Rp_X)
       -\tfrac12R(\mathcal Fp_U-R^{-1}p_U).
\]
Fix the vertex universe to be the union of the original and reflected
spectral labels in $K$, of cardinality $n\le4d$. Retain only cells
$(x,u)\in K^2$ with $\phi(x,u)\le s\le1$, as directed edges $x\to u$.
Keep isolated vertices. If $m_s$ is retained mass and $o,i$ are outgoing
and incoming masses, Markov's inequality gives fixed constants with
\begin{equation}\label{eq:wf-balance}
 1-m_s\le D_0\eps/s,\qquad
 \|i-Ro\|_1\le C_b\sqrt\eps+(1+R_{\max})D_0\eps/s.
\end{equation}
No physical realization of $\bar\mu$ is asserted or needed.

\subsubsection*{Cycles and paths}

An elementary monotonicity observation will be used. If
$x_0,\ldots,x_k=x_0$ satisfies $|x_{j+1}-h(x_j)|\le\delta$ with
$h$ nondecreasing, then $|h(t)-t|\le\delta$ throughout the visited
interval. Every interior level is crossed upwards and downwards; an
upward crossing gives $h(t)\ge t-\delta$, and a downward one gives
$h(t)\le t+\delta$. Predecessor and successor inequalities give the
same bounds at the minimum and maximum. In particular every cycle edge
has length at most $2\delta$.

Put $\Delta=S-1/4>0$. Since $t^2-1+2b(t)\le1/4$, on the diagonal
\[
 \phi(t,t)=S-[t^2-1+b(t)(r(t)+r(t)^{-1})].
\]
For $|R(t)-1|\le c\le1/2$ the identity
$r+r^{-1}-2=(R-1)^2/[r(r+1)^2]$ implies
$\phi(t,t)\ge\Delta-c^2$. Choose
$c=\min(1/2,\sqrt\Delta/2)$. On a retained cycle,
\eqref{eq:wf-predictor} instead bounds
$\phi(x,x)\le s+2L_\phi\delta$, where
$\delta=2\sqrt{Ms}$ and $L_\phi$ is a Lipschitz constant on $K^2$.
Choose a fixed sufficiently small $\delta_*>0$ so that the latter
quantity is at most $\Delta/2$ whenever $\delta\le\delta_*$.
Every cyclic vertex must then have $R\le1-c$ or $R\ge1+c$.
Every internal edge of a cyclic strongly connected component belongs
to a cycle. Decreasing $\delta_*$ so that $2L_R\delta_*<2c$ shows
that such a component is entirely low or entirely high.

Individual components are not enough: a low cycle feeding a high cycle
can support a nonzero exactly balanced flow. To exclude this, also take
$\delta_*\le c/(4\max(1,L_R))$, and choose the cut using the previously
fixed universe:
\begin{equation}\label{eq:wf-cut}
 \delta_n=\delta_*(M+1)^{-n},\qquad s_n=\delta_n^2/(4M).
\end{equation}
Starting at a cyclic vertex $a$, any $k$-edge path satisfies
\[
 |x_{j+1}-a|\le M|x_j-a|+2\delta_n,\qquad
 |x_k-a|\le2\delta_n\sum_{j=0}^{k-1}M^j
              \le2\delta_n(M+1)^k.
\]
A shortest path is simple, hence has at most $n-1$ edges. Its endpoint
cannot be an opposite-gain cyclic vertex: such vertices differ in $R$
by at least $2c$, while the displayed distance and Lipschitz continuity
give a difference at most $c/2$. In particular no low cyclic component
can reach a high cyclic component.

\subsubsection*{A signed potential that controls every branch}

For any $n$-vertex directed graph with these properties there is a real
function $f$ such that
\begin{equation}\label{eq:wf-potential}
 R_xf_x-f_y\ge1\quad(x\to y),\qquad
 \|f\|_\infty\le B_*^{n+1},\quad
 B_*=4+2R_{\max}+2/R_{\min}+2/c.
\end{equation}
Here is the construction. Let $\mathcal U$ be all vertices reachable
from low cyclic components. It is forward closed and contains no high
cycle. Its complement is backward closed, and no strongly connected
component straddles the split: any component meeting $\mathcal U$ is
contained in it by forward closure. On its condensation directed acyclic graph choose constants in
topological order,
\[
 g_C=\max\{1,1/c,\max_{D\to C}(R_{\max}g_D+1)\},
\]
omitting empty inner maxima, and put $f=-g_C$ on $C$. An internal
edge in a cyclic component has $(1-R_x)g_C\ge1$; edges between
components satisfy $g_{C(y)}-R_xg_{C(x)}\ge1$ by construction.
An acyclic singleton has no internal edge. On the complementary
condensation graph process in reverse topological order and set
\[
 f_C=\max\{1/R_{\min},1/c,
       \max_{C\to D,\ D\not\subset\mathcal U}(f_D+1)/R_{\min}\}.
\]
Its cyclic components are high, so internal edges obey
$(R_x-1)f_C\ge1$; edges between components obey the same required
inequality by the recursion. Edges into $\mathcal U$ obey it because
$f_x\ge1/R_{\min}$ and $f_y<0$. No edge leaves $\mathcal U$.
At most $n$ stages, with fixed bounded growth at each, give the norm
bound in \eqref{eq:wf-potential}.

For arbitrary nonnegative weights on all retained edges, summing gives
\begin{equation}\label{eq:wf-all-branches}
 m_s\le\sum_{x\to y}\bar\mu_{xy}(R_xf_x-f_y)
       =\sum_x f_x(R_xo_x-i_x)
       \le B_*^{n+1}\|Ro-i\|_1.
\end{equation}
This is why branching and spectral multiplicity introduce no missing
payment: no edge, row, or component is replaced by a selected path.

\subsubsection*{Assembly}

If the fixed universe is empty, the collar estimate already gives a
positive lower bound on $\eps$. Otherwise combine
\eqref{eq:wf-balance}--\eqref{eq:wf-all-branches}, with
$C_3=(1+R_{\max})D_0$ and constants enlarged to at least one:
\[
 1-D_0\eps/s_n\le B_*^{n+1}
                 (C_b\sqrt\eps+C_3\eps/s_n).
\]
Either $D_0\eps/s_n\ge1/2$, or one of the two right-hand terms is
at least $1/4$. Hence
\[
 \eps\ge\min\left\{\frac{s_n}{2D_0},
 \frac{1}{16C_b^2B_*^{2n+2}},
 \frac{s_n}{4C_3B_*^{n+1}}\right\}.
\]
By \eqref{eq:wf-cut} this is at least $c_0e^{-C_0n}$ with fixed positive
constants. Since $n\le4d$, \eqref{eq:lb} follows. It also excludes exact
finite-dimensional attainment without using spatial attainment or the
old common-return argument. Finally,
$\eps\ge c_*e^{-C_*d}$ implies
$d\ge C_*^{-1}\log(c_*/\eps)$, proving the lower half of
Theorem~\ref{thm:rate}.\hfill$\square$

\paragraph{Verification scope.}
This proof is analytic. Exact finite-graph tests exercise the signed
potential and negative controls, but do not prove the analytic estimates.
The Lean module \texttt{I3322Kernel.WeightedFlow} verifies the finite
all-edge summation identity, the residual bound conditional on an edge
potential, three local potential inequalities, and the elementary
assembly alternative. It does not prove the existence of the potential
from strongly connected components, convex-envelope regularity, or the
quantum-to-flow reduction. The full chain remains an analytic proof.
The earlier staircase and selected-walk cores are not dependencies of
this replacement.
